\chapter{The 2-Categorical Position of Music Among Time-Value Art Forms}

\begingroup

% Local macros copied from the standalone source.

\def\Cat{\mathbf{Cat}}

\def\Set{\mathbf{Set}}

\def\B{\mathbf{B}}

\def\V{\mathcal{V}}

\def\I{\mathcal{I}}

\def\J{\mathcal{J}}

\def\Rhy{\mathbf{Rhy}}

\def\Mel{\mathbf{Mel}}

\def\Harm{\mathbf{Harm}}

\def\Pitch{\mathbf{Pitch}}

\def\Disc{\mathbf{Disc}}

\def\Frame{\mathfrak{A}}

\def\op{^{\mathrm{op}}}

\def\R{\mathbb{R}}

\def\Z{\mathbb{Z}}

\def\N{\mathbb{N}}

\def\Xmus{X_{\mathrm{mus}}}

\begin{abstract}
\noindent This chapter models rhythm, melody, and harmony through a
parameter pair $(\B, \V)$: a base category of temporal or event
structure and a value category of possible labels, pitches, timbres,
or expressive states. Letting this pair vary organises time-value art
forms into a 2-category $\Frame$. The chromatic lattice, viewed as a
finite directed graph, distinguishes a reference object
$\Xmus \in \Frame$. The question \emph{``is $X$ music?''} is
recast as a structural question about the hom-category
$\Frame(X, \Xmus)$: it may admit faithful musicalisations, fully
faithful ones, or none. We classify several borderline cases — paired
poetry, dialogue, Gregorian chant, classical Chinese verse,
\emph{Sprechstimme}, birdsong, ritual percussion, a calling flock —
and read off, in each case, what the verdict depends on.
\end{abstract}

% =====================================================================
\section{From parameters to a 2-category of frames}
% =====================================================================

Rhythm, melody, and harmony can be parametrised by a pair $(\B, \V)$
of small categories. Instead of treating that pair as fixed, we let it
vary.

\begin{definition}[The 2-category of time-value frames]\label{def:frames}
The 2-category $\Frame$ has:
\begin{enumerate}[label=\textup{(\roman*)}, leftmargin=2em]
\item \emph{Objects} (called \emph{frames}): pairs
$X = (\B_X, \V_X)$ of small categories.
\item \emph{1-morphisms} $X \to Y$: pairs
$(\beta, \nu)$ where $\beta : \B_X \to \B_Y$ and $\nu : \V_X \to \V_Y$
are functors.
\item \emph{2-morphisms} $(\beta, \nu) \Rightarrow (\beta', \nu')$:
pairs $(\sigma, \tau)$ of natural transformations
$\sigma : \beta \Rightarrow \beta'$ and $\tau : \nu \Rightarrow \nu'$.
\end{enumerate}
Composition and units are component-wise.
\end{definition}

\begin{construction}[Functoriality of rhythm, melody, harmony]
The constructions of the previous framework extend to 2-functors
$\Frame \to \Cat$:
\begin{align*}
\Rhy   &: X \mapsto \Rhy(\B_X), \\
\Mel   &: X \mapsto \Mel(\B_X, \V_X), \\
\Harm_\J &: X \mapsto \Harm_\J(\B_X, \V_X).
\end{align*}
A 1-morphism $(\beta, \nu) : X \to Y$ acts on rhythms by
post-composition with $\beta$, on melodies by post-composition with
$(\beta, \nu)$ in each leg of the span, and on harmonies
component-wise on each voice.
\end{construction}

\begin{remark}
This is the meta-theoretic move: every art form whose phenomenology
is ``events in some base, labelled by some values'' is now an object
of $\Frame$, and the rhythm/melody/harmony machinery applies to all
of them uniformly. What distinguishes \emph{music} from its
neighbours is not the machinery but the choice of frame.
\end{remark}

% =====================================================================
\section{The chromatic graph}
% =====================================================================

The reference frame for music is built from a graph, not just a set.

\begin{definition}[Chromatic graph]\label{def:chromatic}
The \emph{chromatic graph} $\Gamma$ has vertex set $\Z$ and a single
directed edge $n \to n{+}1$ for each $n \in \Z$. Its \emph{free
category} is the chromatic lattice
\[
\Pitch \;:=\; \mathrm{Free}(\Gamma).
\]
Concretely, $\mathrm{Hom}_\Pitch(m, n) = \{*\}$ if $n \geq m$ and
$\emptyset$ otherwise; composition is the unique map. Equivalently,
$\Pitch \cong (\Z, \leq)$ as a thin category.
\end{definition}

\begin{remark}[Why graph and not just set]
$\Pitch$ as a discrete set $\Disc(\Z)$ records pitch identity only.
$\Pitch$ as the free category on $\Gamma$ additionally records
\emph{intervallic order}: the unique morphism $m \to n$ exists iff $n$
lies above $m$, and the length of the shortest factorisation is
$n - m$ semitones. Functors $\V_X \to \Pitch$ must therefore preserve
whatever ordered structure $\V_X$ carries, not merely the underlying
set of values. This distinction is the \emph{entire content} of the
question that follows.
\end{remark}

\begin{definition}[Finite ambitus, groupoidification]
For $a \leq b$ in $\Z$, $\Pitch_{[a,b]}$ is the full subcategory on
$\{a, a{+}1, \ldots, b\}$. The \emph{interval groupoid}
$\Pitch^\sim$ is the localisation of $\Pitch$ at all morphisms; it
has the same objects with $\mathrm{Hom}_{\Pitch^\sim}(m, n) = \{*\}$
for all $m, n$. Inversion symmetry — descending intervals as
inverses of ascending ones — lives in $\Pitch^\sim$.
\end{definition}

\begin{definition}[Musical reference frame]
The \emph{musical reference frame} is
\[
\Xmus \;:=\; \bigl((\R_{\geq 0}, \leq),\ \Pitch\bigr) \;\in\; \Frame.
\]
\end{definition}

% =====================================================================
\section{Musicalisation and the musical predicate}
% =====================================================================

\begin{definition}[Musicalisation]
A \emph{musicalisation} of a frame $X \in \Frame$ is a 1-morphism
\[
\Phi = (\beta, \nu) : X \longrightarrow \Xmus
\]
in $\Frame$. The \emph{musicalisation category} of $X$ is the
hom-category $\Frame(X, \Xmus)$, whose objects are musicalisations
and whose morphisms are 2-cells.
\end{definition}

\begin{definition}[The musical predicate]\label{def:predicate}
A frame $X$ is:
\begin{itemize}[leftmargin=2em]
\item \emph{transcribable} if $\Frame(X, \Xmus)$ is non-empty;
\item \emph{musical} if it contains a 1-morphism $(\beta, \nu)$ with
both $\beta$ and $\nu$ faithful;
\item \emph{strictly musical} if it contains one with both
$\beta$ and $\nu$ fully faithful;
\item \emph{canonically musical} if the connected component of any
faithful $\Phi$ is unique up to a chosen 2-isomorphism class;
\item \emph{amusical} if $\Frame(X, \Xmus) = \emptyset$.
\end{itemize}
\end{definition}

\begin{remark}[The four-step ladder]
The conditions above form a ladder:
\[
\text{amusical}
\;<\; \text{transcribable}
\;<\; \text{musical}
\;<\; \text{strictly musical},
\]
with \emph{canonically musical} as a transverse property concerning
the \emph{number} of faithful musicalisations rather than their
existence. Most everyday verdicts about ``what is music'' confuse
these levels; the framework lets us separate them.
\end{remark}

% =====================================================================
\section{Worked examples}
% =====================================================================

In every example below, the time-leg $\beta$ is the identity
inclusion $(\R_{\geq 0}, \leq) \hookrightarrow (\R_{\geq 0}, \leq)$.
The verdict is therefore controlled by the value-leg
$\nu : \V_X \to \Pitch$.

% --------------------------------
\subsection{Two poets in alternation}
% --------------------------------
Let $\V_{\mathrm{phon}}^\flat := \Disc(P)$, the discrete category on
the finite set $P$ of phonemes of the poets' language. Any injection
$\nu : P \hookrightarrow \Z$ — alphabetical, frequential, arbitrary
— gives a fully faithful functor $\V_{\mathrm{phon}}^\flat \to \Pitch$,
because both source and target are discrete on their objects.

\begin{verdict}
Two poets are \emph{strictly musical} but \emph{not canonically
musical}: the musicalisation exists, even fully faithfully, but no
choice of $\nu$ is privileged.
\end{verdict}

\noindent The framework here corrects a common intuition: the
familiar judgement \emph{``poetry is not music''} does not rest on a
structural failure but on the absence of a \emph{distinguished}
musicalisation. Cultures that fix one (e.g.\ liturgical psalmody, Vedic
recitation) immediately upgrade the same raw structure to canonically
musical.

% --------------------------------
\subsection{Conversation}
% --------------------------------
Now enrich the value category to $\V_{\mathrm{conv}}$: same objects $P$,
but with non-identity morphisms encoding minimal-pair relations
($/p/ \to /b/$ for voicing, $/n/ \to /m/$ for place, etc.). A functor
$\nu : \V_{\mathrm{conv}} \to \Pitch$ must send each minimal-pair
morphism to a chromatic morphism, i.e., must respect a partial order
on $P$ compatible with the chromatic order on $\Z$. The minimal-pair
graph is in general not a sub-poset of $(\Z, \leq)$ — it has cycles,
incomparable pairs, and no consistent direction.

\begin{verdict}
Conversation is \emph{transcribable} (drop the morphisms, fall back
to the discrete case) but \emph{not musical}: no faithful $\nu$
respects the linguistic structure.
\end{verdict}

\noindent The boundary between this and the poets case is exactly
how richly one models $\V$: stripping the value category to its
discrete shadow is the act of treating speech as music.

% --------------------------------
\subsection{Gregorian chant}
% --------------------------------
$\V_{\mathrm{chant}} := \Pitch_{[a,b]}$ for some Dorian or
Mixolydian ambitus. The inclusion
$\V_{\mathrm{chant}} \hookrightarrow \Pitch$ is fully faithful and
even essentially surjective onto its image.

\begin{verdict}
Gregorian chant is \emph{strictly musical} and \emph{canonically
musical}: the musicalisation is the inclusion, with no choice
involved.
\end{verdict}

% --------------------------------
\subsection{Classical Chinese verse}
% --------------------------------
$\V_{\mathrm{verse}} := \V_{\mathrm{phon}}^\flat \times \Disc(\Z/4)$,
where the second factor records lexical tone (level, rising,
departing, entering). Two natural musicalisations present themselves:
\begin{align*}
\nu_{\mathrm{tone}} &: \V_{\mathrm{verse}}
   \xrightarrow{\mathrm{pr}_2} \Disc(\Z/4)
   \hookrightarrow \Pitch / 12\Z, \\
\nu_{\mathrm{full}} &: \V_{\mathrm{verse}}
   \xrightarrow{\mathrm{any\ injection}} \Pitch.
\end{align*}
The first is canonical but not faithful (it forgets phonemes). The
second is faithful but not canonical (any injection works).

\begin{verdict}
Classical Chinese verse is \emph{musical} and admits a
\emph{canonical quotient-musicalisation} via tone, but not a
canonical full musicalisation. The genre's lived ambiguity between
``poetry'' and ``song'' is exactly this split.
\end{verdict}

% --------------------------------
\subsection{Sprechstimme}
% --------------------------------
$\V_{\mathrm{spr}} := \Pitch \times \Disc(\{\mathrm{spoken}\})$. The
first projection $\mathrm{pr}_1 : \V_{\mathrm{spr}} \to \Pitch$ is
fully faithful (the second factor is terminal).

\begin{verdict}
\emph{Sprechstimme} is \emph{strictly musical} and \emph{canonically
musical}. The \texttt{spoken} marker is a stylistic decoration that
the framework correctly recognises as inessential to the
classification.
\end{verdict}

% --------------------------------
\subsection{Birdsong}
% --------------------------------
$\V_{\mathrm{bird}} := (\R, \leq)$, continuous pitch contour with its
order. A faithful functor $(\R, \leq) \to (\Z, \leq) \cong \Pitch$
would require an order-injection of an uncountable set into a
countable one — impossible.

\begin{verdict}
Birdsong is \emph{transcribable} (the rounding map
$\R \to \Z$ is a non-faithful musicalisation) but
\emph{not musical}. Western notation of birdsong is precisely this
quantisation, and the loss of information is structural, not
artisanal.
\end{verdict}

% --------------------------------
\subsection{Ritual percussion of arbitrary timbres}
% --------------------------------
$\V_{\mathrm{drum}} := \Disc(T)$ for a finite set $T$ of timbres.
This is the two-poets situation again: any injection $T \to \Z$ is a
fully faithful musicalisation; none is canonical.

\begin{verdict}
Strictly musical, not canonically musical. The standard practice of
assigning timbres to lines of a percussion staff is the social act of
fixing $\nu$ by convention.
\end{verdict}

% --------------------------------
\subsection{A flock of birds calling simultaneously}
% --------------------------------
This is a harmony in $\Harm_\J(\B, \V_{\mathrm{bird}})$ for
$|\J|$ = the number of birds. Musicalisability is inherited from the
voice components: a harmony is (strictly) musical iff each voice is.

\begin{verdict}
Amusical, on the same grounds as a single bird. The framework
\emph{cannot} confer musicality on an ensemble whose components lack
it.
\end{verdict}

% =====================================================================
\section{Summary table}
% =====================================================================

\begin{center}\small
\renewcommand{\arraystretch}{1.25}
\begin{tabular}{lll}
\hline
\textbf{Frame} & \textbf{Verdict} & \textbf{Decided by} \\
\hline
Two poets               & strictly, not canonically  & freedom of $\nu$ \\
Conversation            & transcribable, not musical & morphisms in $\V_{\mathrm{conv}}$ \\
Gregorian chant         & strictly + canonically     & inclusion of ambitus \\
Classical Chinese verse & quotient-canonical         & tone projection \\
Sprechstimme            & strictly + canonically     & terminal second factor \\
Birdsong                & transcribable, not musical & cardinality \\
Ritual percussion       & strictly, not canonically  & freedom of $\nu$ \\
Calling flock           & amusical                   & inherited \\
\hline
\end{tabular}
\end{center}

% =====================================================================
\section{Closing: what the framework refuses to decide}
% =====================================================================

The framework converts \emph{``is $X$ music?''} into a precise
question about which 1-morphisms exist in $\Frame(X, \Xmus)$, and it
exposes the variables on which a verdict depends:
\begin{enumerate}[label=\textup{(\roman*)}, leftmargin=2em]
\item how richly one models $\V_X$ (discrete shadow vs.\ full
morphism structure);
\item whether one demands the musicalisation to be canonical or
merely possible;
\item whether one accepts $\Xmus$ as fixed or allows alternative
reference frames (microtonal $\Pitch_{\mathrm{micro}}$, just-intonation
$\Pitch_{\mathrm{ji}}$, electroacoustic spectrum
$\Pitch_{\mathrm{spec}}$, etc.).
\end{enumerate}

\noindent What the framework \emph{does not} decide is which of these
choices is the right one: the choice is a substantive musicological
or philosophical commitment, and different communities make it
differently. The framework's contribution is to make the disagreement
\emph{locatable}: two interlocutors who reach opposite verdicts on
``poetry as music'' will, on inspection, disagree about whether to
take $\V_{\mathrm{phon}}^\flat$ or $\V_{\mathrm{conv}}$, and that
disagreement has a name.

\begin{remark}[Why this is a useful refactor]
A philosophical question of the form \emph{``is $X$ really $Y$?''} is
typically irresolvable as posed because $Y$ is given by ostension
rather than by definition. The categorical reframing replaces $Y$ by
a structured object ($\Xmus$) and the predicate ``really $Y$'' by a
class of structure-preserving maps into it. Whether or not one
accepts the chromatic frame as the right reference for ``music'',
one at least now knows what one is rejecting.
\end{remark}

\endgroup
